% Draft of the new introduction, EV-first framing.
% Workspace file — not yet wired into main.tex. Mirrors the chapter scaffolding
% of 00-introduction.tex (\chapter, \label, per-chapter \printbibliography)
% so that switching it in is a one-line change in main.tex.
%
% §1 is the only section drafted here. Companion notes:
%   thesis/notes/intro/outline.md
%   thesis/notes/intro/section-1-outline.md
%   thesis/notes/intro/section-1-citations.md

\chapter{Introduction}
\label{ch:introduction}

\section{The external-validity gap in empirical economics}
\label{sec:intro:ev-gap}

External validity is the binding open problem of empirical economics, and it is field-internal critics --- \citeauthor{Manski2013}, \citeauthor{Heckman2008}, \citeauthor{Deaton2010}, and, from inside the randomised-controlled-trial movement itself, \citeauthor{Snowberg2016} --- who have pressed the point for decades. \citet{Manski2013} writes that ``analyses of experimental data have tended to be silent on the problem of extrapolating from the experiments performed to policies of interest'' (p.~37); the credibility of a study, in his account, ``depends jointly on internal and external validity'' (pp.~36--37), so privileging the first at the cost of the second leaves the policy question unanswered. \citet{Heckman2008} sharpens the diagnosis from the structural side: forecasting ``the effects of interventions in environments other than the one where the estimate was obtained'' --- ``the problem of external validity'' (p.~8) --- is a task that randomised estimation, by construction, does not deliver. \citet{Deaton2018} restate the critique in its strongest contemporary form: randomisation buys internal validity by design, but ``the lack of structure is often a disadvantage when we try to use the results outside of the context in which the results were obtained; credibility in estimation can lead to incredibility in use'' (p.~3); see also \citet{Deaton2010}. We use \emph{external validity} throughout in the sense of \citet[p.~5]{CampbellStanley1963}, who introduced the internal/external distinction and named the question it asks --- ``to what populations, settings, treatment variables, and measurement variables can this effect be generalized?'' --- and who endorsed the primacy-of-internal-validity ordering that the econometric critics later rejected. The lineage's punctuation comes from inside the RCT movement itself: \citet{Snowberg2016} write that ``creating a rigorous framework for external validity is an important step in completing an ecosystem for social science field experiments.'' That EV is the unfinished business of empirical economics is, in other words, not an outsider critique. It is a recognition --- not yet an operational fix.

The recognition is not theoretical. A body of empirical work demonstrates that randomised estimates often fail to replicate across contexts, populations, and time periods --- by enough, in places, to invert the standard hierarchy of evidence. \citet{Pritchett2015} run a direct horse race on six microcredit experiments and find that within-context observational estimates produce a smaller root-mean-squared prediction error than between-context experimental estimates: 0.18 against 0.33 for business profits, 0.36 against 0.49 for household consumption (p.~473). \citet{Allcott2015}, studying the Opower Home Energy Reports program across 111 utilities, shows that the early experimental sites overpredicted the at-scale average treatment effect by 0.41--0.66 percentage points --- a documented \emph{site-selection bias} in which utilities that adopt the program early have systematically more treatment-responsive customers (p.~1119). \citet{Bisbee2017} document an analogous failure for instrumental-variable estimates: across 139 country-year combinations of the same-sex-instrument LATE for the effect of fertility on female labour supply, cross-context extrapolation of a clean local average treatment effect performs systematically worse than within-context observational estimation. \citet{Rosenzweig2019} add the temporal axis: in panel data on returns to agricultural investment, the realisation of aggregate shocks is large enough that the probability that any single year's estimate lies within thirty percentage points of the expected rate of return is only five per cent in both Ghana and India. The site-selection bias of \citet{Allcott2015} and the complier heterogeneity of \citet{Bisbee2017} are, at root, \emph{representativity} failures: the experimental sample is not drawn from the population the policy will act on (Allcott), and the compliers driving the local average treatment effect are not the compliers a policy would induce in a different context (Bisbee). Representativity is not a survey-side concern adjacent to RCT methodology; it is a documented driver of RCT external-validity failure in the empirical record itself.

External validity is therefore not a single property of an estimate but a bundle: whether a causal effect transports depends jointly on the population it was estimated on, the setting in which it was estimated, the treatment as it was operationalised, and the outcome as it was measured \citep{CampbellStanley1963}. Two of these sub-dimensions have crystallised into terms of art in their own right. By \emph{representativity} we mean the persons-dimension property that the sample's joint distribution on outcome-relevant covariates matches that of the target population --- the standard sense from the survey-sampling tradition of \citet{Neyman1934}, and the property whose empirical failure was documented by \citet{Allcott2015} and \citet{Bisbee2017}. By \emph{ecological validity} we mean the bundle of the remaining three sub-dimensions on the experimental-design side: the resemblance of the experimental setting, the stimulus as delivered, and the response as measured to the conditions an intervention would meet were it actually deployed --- the three-dimensional formulation of \citet{Schmuckler2001}. Experimental economics has separately characterised the same property through the six field-context dimensions of \citet{Harrison2004} --- subject pool, information, commodity, stakes, task, and environment --- whose joint match to the policy-relevant context defines a \emph{natural field experiment}; we adopt Schmuckler's noun and Harrison--List's operational dimensions in tandem. A further failure mode is now well documented: pilot effects shrinking or vanishing at scale, the \emph{voltage drop} of \citet{Al-Ubaydli2020}, with empirical demonstration in \citet{Bold2018} and programmatic statement in \citet{Muralidharan2017}. We treat this not as a third umbrella term but as a special case of ecological validity at the largest setting scale, and address it jointly with ecological validity in Chapter~\ref{ch:malaria}. This thesis develops methods for two of these factors --- representativity in Chapter~\ref{ch:survey-sampling} and ecological validity (inclusive of the at-scale aspect) in Chapter~\ref{ch:malaria}. Other factors named in the prior literature --- complier composition, mechanism invariance, construct stability --- are real and not addressed here. The two factors taken up are not symmetric: representativity is what a pure-survey design pursues directly and what an experimental design needs as a precondition for the intention-to-treat effect to have a meaningful target population, while ecological validity is an experimental-design property with no pure-survey analogue. The next section makes that asymmetry precise.

\printbibliography[heading=subbibliography, title={References for Chapter \thechapter}]
